\section{Conclusion}
We presented the radiometric camera, which lifts any central geometric
camera model to a radiometric one with a finite aperture, a free-form
focus surface and a generic vignetting function. We also presented its
calibration from the RAW photosites of ordinary calibration images,
formulated without parameter-dependent discontinuities or integration
domains and with every gauge freedom fixed exactly. To our knowledge,
it is the first model to unite aperture, focus and vignetting with
any central projection and to calibrate all of them jointly from the
images of a geometric calibration (\cref{tab:related}). The calibrated
cameras have a lower reprojection error on a held-out high-precision
target than their geometric starts for most models and lenses, their
vignetting reproduces a measured flat field, and their focus surface
is learned as a plane for a wide-angle lens and as a hemisphere for a
fisheye. In a first multi-view radiometric reconstruction, they
recover the in-focus reflectance of a texture from defocused images.
The estimation of light sources and of the attenuation of water
is future work. We will
release the code, the data and the configurations of every experiment.

% ---- Ties's conclusion of 11-17, kept for reference ----
% A generative camera model enables the learning of a generative scene.
% A bias in the geometric model as well as a bias in the radiometric model will propagate into the scene.
% Creating an accurate geometric camera model for an unbiased 3D reconstruction has been tried for several years.
% However, even the most accurate geometric model cannot be used to reconstruct scene radiometry since the geometric model cannot match a radiometric sensor.
% We showed that our camera model is indeed a generative model: we can learn geometric distortion models combined with arbitrary focus surfaces and arbitrary vignetting, matching the physical properties of our calibrated lenses.
% Our generative model gave us already the possibility to reach a smaller reprojection error in calibration, a joint learning of the vignetting and a sharp texture reconstruction from out-of-focus images.
% There are also challenges arising from radiometric optimization:
% we had to cope missing constrains and several gauge freedoms for our calibration.
% Moreover, differentiating integrals during differentiable ray tracing entails the risk of missing boundary terms leading to wrong gradients.
